\section{Introduction}\label{sec:introduction}
Consider a sum of functions, each depending on a small subset of the
variables. A tree decomposition can separate most of the variables into
small overlapping bags. It is natural to use this structure in a global
optimization proof: bound each bag locally, communicate lower bounds on
the overlaps, and combine the bounds along the tree. The difficulty is
that the communicated functions are generally nonconvex, nonsmooth, or
expensive to represent. Replacing them by constants loses first-order
information. Replacing them by affine functions on small cells preserves
more information, but raises a different question: how many cells are
needed, and how can an algorithm find them?

This paper studies a finite certificate with three ingredients: convex
lower models on bag boxes, affine minorants on separator cells, and local
inequalities that can be checked independently. A lower certificate and
a feasible upper value certify an absolute objective gap. Our main result
constructs these objects by \emph{regridding}: at each stage, it discards
the previous partitions and builds fresh graded partitions around a
consistent point recovered from the previous lower-bound calculation.
The new partitions are finer near that point and coarser farther away.
Their separator slopes are subtree sums of bag gradients evaluated at
the same point.

The assumptions are substantive. The objective has global quadratic
growth about a unique minimizer, each bag gradient is Lipschitz, and each
local lower model has an error bounded by a constant times the square of
the box width. The structural parameters are the maximum bag dimension
and the maximum number of bags containing a coordinate. The constants
also depend on the ratios of curvature and model error to quadratic
growth. They do not depend on the branching degree or the depth of the
decomposition. The algorithm does not know the minimizer and does not
use stationarity at it. Boundary minimizers are included.

\subsection{Contributions and computational guarantees}
The principal contribution is a construction bound on arbitrary tree
decompositions. A configuration attaining the lower bound may assign
different values to copies of the same coordinate. Those disagreements
cannot in general be controlled by the largest distance of a bag from the
current center. We instead bound their total squared norm. This gives a
contraction of the recovered consistent point toward the minimizer and
also bounds its objective gap. For a sufficient fixed grading ratio, the
number of final boxes and cells is linear in the number $N$ of bags and
logarithmic in accuracy, with an exponential dependence on bag dimension.

The distinction between a certificate and its construction is essential.
Let $p$ be the maximum bag dimension, $q$ the maximum separator dimension,
and $\theta$ the grading ratio. Let $J$ denote the number of scales needed
to reach the requested accuracy. The explicit definition of $J$ appears
in \cref{thm:regridding-main}; for fixed instance constants,
$J=O(1+\log_+(N/\eps))$. The main quantities are
\begin{center}
\begin{tabular}{@{}ll@{}}
\toprule
Quantity & Bound for a sufficient fixed $\theta$\\
\midrule
Final boxes and cells & $O\bigl(N(4/\theta)^p(J+1)\bigr)$\\
Boxes and cells created over all stages & $O\bigl(N(4/\theta)^p(J+1)^2\bigr)$\\
Local convex minimizations & $O\bigl(N3^q(4/\theta)^p(J+1)^3\bigr)$\\
\bottomrule
\end{tabular}
\end{center}
These are counts of specified objects and oracle calls. They do not bound
the internal work of an arbitrary convex-minimization oracle. The paper
provides two ways to make this distinction operational. First, certified
inexact solves are allowed: local errors are summed along a selected
configuration rather than over every subproblem. Second, rational
polynomial objectives of fixed degree admit affine Taylor lower models,
for which the local problems reduce to exact endpoint minimization.
This gives explicit bit and output bounds. The resulting dependence on
conditioning is numerical, so the theorem is not a polynomial-time claim
for arbitrary rational input with unbounded conditioning.

The output is a global lower proof together with a feasible upper value,
rather than only an approximate solution. In the polynomial realization,
an independent checker can reconstruct every local affine model and verify
the gap by rational arithmetic. It does not need to verify the growth
constant used to analyze construction time. This separation is useful
when a computed optimum must be checked independently, or when a sparse
problem must be solved to high precision with an explicit proof-size bound.

A budgeted search removes the need to know the sufficient grading ratio.
Every trial remains sound, and a successful sufficient trial gives a
bound on the total search work. A trial that stops earlier may use a
different ratio and a different scale. Its validity and output budget are
guaranteed; the sharper final partition bound above belongs to the
sufficient fixed-ratio run. We keep these statements separate.

The supporting theory has two purposes. We first quantify what is lost
when different bags use inconsistent copies of a separator coordinate.
For one separator, the best gap in a permitted linear class of shifts
equals twice the uniform distance from that class to the admissible band
between the child value function and the opposite-side lower requirement.
Independent cell pieces reduce the gap to the largest cell bracket.
Tree versions give lower and upper bounds in
terms of the local deficits. The available conditional margin discounts
approximation error away from optimal separator values. This explains
why uniform approximation of every conditional value function can be
unnecessarily costly. A scalar covering result makes this observation
quantitative.

We then extend the construction to scalar contracting nonlinear dynamics.
Box partitions are intersected with affine strips containing the local
graphs. Forward simulation repairs a relaxed configuration into an
exactly feasible trajectory. Ordinary objective-gradient slopes are not
enough under coupling constraints; we use slopes adjusted by a backward
adjoint recursion. This recursion cancels state derivatives at arbitrary
ambient centers, which also permits finite-precision resets. In rational
polynomial dynamics, exact expansion of a trajectory can have very large
bit length. A recurrence representation, together with certified
enclosures and an objective bound, supplies a smaller exactly feasible
output.

\subsection{Relation to prior work}\label{sec:prior-work}
Separator messages, cost shifting, and optimization over a tree are
established ideas. In finite-state graphical models, agreement on trees
leads to linear-programming and reparameterization interpretations
\citep{WainwrightEtAl2005}; AND/OR search exploits separator-conditioned
subproblems and cached states \citep{MarinescuDechter2009}.
Weighted constraint satisfaction combines tree decomposition with local
consistency and optimized cost shifts
\citep{deGivrySchiexVerfaillie2006,CooperDeGivrySchiex2007}.
For continuous pairwise models, Wald and Globerson
\citep{WaldGloberson2014} impose agreement on selected expectations and
derive a dual consisting of cost-shifted local minima. Their exactness
result for convex-decomposable pairwise models has a different hypothesis
from our growth-based construction. The restricted-split duality in
\cref{thm:consistency-duality} belongs to this line of work; we give its
compact minimax proof for the function classes used here. Unrestricted
tree gluing is also an established mechanism in sparse polynomial
optimization \citep{Lasserre2006}.

Positive-margin splitting under running intersection is central to sparse
polynomial representations. Grimm, Netzer, and Schweighofer
\citep{GrimmNetzerSchweighofer2007} construct such splits qualitatively;
Korda, Magron, and R\'ios-Zertuche
\citep{KordaMagronRiosZertuche2025} quantify polynomial approximation of
separator functions. Local polynomial shifts also appear in tightness
criteria for sparse moment relaxations \citep{NieQuTangZhang2026}.
Factor-two approximation bounds have precedents in approximate linear
programming \citep{deFariasVanRoy2003}, and moment matching gives a
zero-width approximation duality \citep{HanJiaoWeissman2018}.
Approximation of interval-valued bands through oriented distance is
studied by Ginchev and Hoffmann \citep{GinchevHoffmann2002}.
For a real interval band, their objective is
$\sup_s\max\{L(s)-\chi(s),\chi(s)-U(s)\}$. When the band pinches,
this is the uniform distance from $\chi$ to the band. Thus this
approximation objective is established.
We do not claim band approximation, qualitative splitting, or the
factor-two approximation principle as new. Our band results identify the
further equality between the optimized restricted dynamic-programming
gap and twice this distance, and distinguish separate-edge approximation
from one jointly exact tree split. The graded split and scalar covering
proof use the margin left in these bands.

Continuous decomposition methods provide closer algorithmic comparisons.
Berenguel et al. \citep{BerenguelEtAl2013} study interval branch-and-bound
for additively separable functions with common variables. Their
shared-variable pruning rule combines interval lower bounds from retained
subboxes whose common-variable intervals overlap the current interval
\citep[Proposition~1 and Theorem~2]{BerenguelEtAl2013}.
This supplies a direct antecedent for compatible-box decomposition bounds.
Our construction communicates affine separator minorants on a tree and
bounds the size and work of the resulting global certificate.
Bienstock and Mu\~noz \citep{BienstockMunoz2018} construct
bounded-treewidth linear-programming formulations for polynomial
optimization, with explicit dependence on inverse accuracy.
Zhang and Sun \citep{ZhangSun2022} give accuracy-dependent iteration
bounds for multistage nonconvex decomposition under their state-covering
and oracle assumptions. MUSE-BB \citep{LangiuEtAl2025} uses decomposable
scenario subproblems and multisection branching in one spatial search
tree. These methods establish the value of decomposition, but their
stated guarantees differ from our final affine-certificate size and
construction count under point growth.
Robertson, Cheng, and Scott \citep{RobertsonChengScott2025} analyze
convergence orders of reduced-space value-function relaxations for
nonconvex stochastic programs.

Adaptive grids are likewise established. Munos and Moore
\citep{MunosMoore2002} refine state-space cells and rebuild discretized
dynamic programs in optimal control. Nagarajan et al.
\citep{NagarajanEtAl2019} combine adaptive multivariate partitioning,
piecewise relaxations, and bound tightening for global optimization.
Our contribution is not adaptivity itself. It is the aggregate disagreement
estimate that makes a fresh graded partition contract around a point
obtained from an affine separator certificate, on an arbitrary supplied
tree. The logarithmic accuracy dependence relies on global quadratic
growth and the displayed structural and conditioning parameters.
Logarithmic accuracy dependence also occurs in optimistic optimization
under suitable local-growth and partition assumptions
\citep{Munos2011OptimisticOptimization}. That guarantee concerns simple
regret. Certified Lipschitz optimization has instance-dependent value-oracle
bounds \citep{BachocCesariGerchinovitz2021}. Our certificate and work counts
use different information: bag gradients, local lower models, and local
convex minimizations. These oracle models should not be compared as if
they supplied the same operations.

Inexact dynamic programming and inexact first-order methods have their own
error-propagation theory \citep{Munos2005,DevolderGlineurNesterov2014}.
Guigues develops deterministic and stochastic inexact decomposition from
approximate primal and dual subproblem solves
\citep{Guigues2017InexactCuts,Guigues2020InexactSDDP}. These related analyses
track stage errors and permit increasingly accurate solves; they already
establish inexact decomposition as a method.
Reliable rational lower bounds from numerical linear programming are
developed by Steffy and Wolter \citep{SteffyWolter2013}, and
limited-precision LP oracles can be used to recover exact solutions
\citep{GleixnerSteffy2020}. We use a specific certified contract: an exact
local feasible point, a valid lower bound, an error budget proportional to
squared cell width, and an aggregate separator-slope error. Its residual
identity charges one local error per selected bag. For polynomial box
objectives, a separate affine Taylor construction eliminates the general
convex oracle: every local affine problem is solved at interval endpoints,
and the resulting rational certificate has a compact representation.
For fixed degree, bag dimension, occurrence, and numerical conditioning,
the bit cost is polynomial in input length and precision. It is not
uniformly polynomial in the curvature-to-growth ratio. This realization
selects a particular aggregate bag model; it does not implement every
possible factorwise relaxation.

Backward adjoints are standard in nonlinear optimal control
\citep{Hager2000}. The dynamics extension uses them to cancel the
first-order objective error of forward feasibility repair in a separator
certificate. The cancellation holds at ambient-box centers as well as
feasible centers. Contracting graph maps then permit short certified
enclosures and repeated center resets. The assumptions include full-box
invariance, feasible point growth, and sound derivative bounds; this is a
specific constrained extension rather than a theorem for arbitrary
nonlinear constraints.

To the best of our knowledge, the results reviewed above do not give the
combination proved here: construction of affine separator certificates
without a known optimizer on arbitrary tree decompositions, the final
logarithmic partition count and explicit incidence-based oracle count
under point growth, certified local error budgets for that construction,
and the stated rational implementations. The claims concern this
combination and these precise contracts. They do not assert novelty of
tree dynamic programming, cost shifting, adaptive partitioning, adjoints,
or exact rational linear programming.


\subsection{Organization}
\Cref{sec:model} defines the problem, the certificate, and its dynamic
program. \Cref{sec:consistency} studies separator consistency and its
approximation cost. \Cref{sec:regridding} proves the main construction
bound; \cref{sec:inexact,sec:arithmetic} treat certified errors and rational
arithmetic. \Cref{sec:dynamics,sec:dynamics-bit} give the nonlinear dynamics
extension and its output representation. \Cref{sec:limitations} proves
examples that delimit the assumptions. The appendices give complementary
covering results with their additional hypotheses stated explicitly.
